\begin{problem}{Self-Crossing Turtle}{standard input}{standard output}{2 seconds}{256 megabytes}

A turtle starts at point $(0, 0)$ on the $X$-$Y$ plane. You are given a sequence of $n$ positive integers $d_1, d_2, \ldots, d_n$.

The turtle moves $d_1$ meters north, then $d_2$ meters west, $d_3$ meters south, $d_4$ meters east, and so on. After every move, its direction rotates by $90^\circ$ counter-clockwise.

Determine whether the path drawn by the turtle intersects itself. Two path segments are considered intersecting even if they only touch at one point or overlap along a non-empty segment. The common endpoint of two consecutive moves is not considered a self-intersection.


\InputFile
The first line contains an integer $t$ ($1 \le t \le 10^4$) --- the number of test cases.

Each test case consists of two lines. The first line contains an integer $n$ ($1 \le n \le 2 \cdot 10^5$) --- the number of moves. The second line contains $n$ integers $d_1, d_2, \ldots, d_n$ ($1 \le d_i \le 10^9$) --- the corresponding distances.

It is guaranteed that the sum of $n$ over all test cases does not exceed $2 \cdot 10^5$.


\OutputFile
For each test case, print $\texttt{Yes}$ if the turtle's path intersects itself, and $\texttt{No}$ otherwise. The answer is case-insensitive: for example, $\texttt{Yes}$, $\texttt{YES}$ and $\texttt{yEs}$ are all accepted.


\Example

\begin{example}
\exmpfile{example.01}{example.01.a}%
\end{example}

\Note
In the first test case, the fourth segment intersects the first segment at $(0,1)$.

In the second test case, no two non-consecutive segments intersect.

In the third test case, the fourth segment returns to $(0,0)$, which belongs to the first segment.

\end{problem}

